% GENERATED FILE. Edit canonical lessons, then run npm run generate:books.
% canonical-source-hash: fnv1a64:cb190962b9a76860
% canonical-lessons: GU-C355-vero, GU-C355-phalo, GU-C355-arthik, GU-C355-grant, GU-C355-ughravvu

\chapter{Taxes and Public Money}
\label{ch:gu-taxes-and-public-money}
\hlchaptermodality{355}

\begin{chapteropening}
\textbf{By the end of this chapter:} \emph{I can talk about taxes and where public money goes.}

Complete the last lesson of chapter 355: I can talk about taxes and where public money goes.
\end{chapteropening}

\section[\texorpdfstring{\gu{વેરો} (vero)}{વેરો (vero)}]{\gu{વેરો} (vero) — a tax}

\label{lesson:GU-C355-vero}



\begin{quote}
Before the new one: say the Gujarati for permission, then the Gujarati for to appoint.
\end{quote}

\subsection*{You'll want to know: \gu{વેરો}}
\textbf{\gu{વેરો}} — \emph{vero} — \textquotedblleft{}a tax\textquotedblright{}.

\begin{quote}
\gu{દરેક} \gu{વરસે} \gu{અમારે} \gu{ઘરનો} \gu{વેરો} \gu{ભરવો} \gu{પડે} \gu{છે}.
\end{quote}

\begin{quote}
\emph{darek varse amāre gharno vero bharvo paḍe che.}
\end{quote}

\begin{quote}
Every year we have to pay the house tax.
\end{quote}

A masculine noun.

\begin{cousinweb}
\textbf{\gu{વેરો}} (\emph{vero}) No dictionary we use records an older form, so its history is not settled.
\end{cousinweb}

\begin{grammarlens}[title={What you've built}]
One more word for giving reasons and opinions.
\end{grammarlens}

\subsection*{Your turn}
\begin{itemize}
\raggedright
  \item \emph{Say it:} \emph{vero}
  \item \emph{Say it:} \emph{vero}, once more
  \item \emph{Say it:} \emph{darek varse amāre gharno vero bharvo paḍe che.} — Every year we have to pay the house tax.
  \item \emph{Recall:} say the Gujarati for permission, then the Gujarati for to appoint, then say \emph{vero} again
\end{itemize}

\begin{tcolorbox}[breakable,colback=teal!4,colframe=teal!35!black,title={Before you move on}]
What does \gu{વેરો} mean? (\textquotedblleft{}A tax\textquotedblright{}.) Say it once more.
\end{tcolorbox}
\section[\texorpdfstring{\gu{ફાળો} (phāḷo)}{ફાળો (phāḷo)}]{\gu{ફાળો} (phāḷo) — a contribution}

\label{lesson:GU-C355-phalo}



\begin{quote}
Before the new one: say the Gujarati for to appoint, then the Gujarati for a tax.
\end{quote}

\subsection*{You'll want to know: \gu{ફાળો}}
\textbf{\gu{ફાળો}} — \emph{phāḷo} — \textquotedblleft{}a contribution\textquotedblright{}.

\begin{quote}
\gu{મંદિર} \gu{માટે} \gu{દરેક} \gu{કુટુંબે} \gu{ફાળો} \gu{આપ્યો}.
\end{quote}

\begin{quote}
\emph{mandir māṭe darek kuṭumbe phāḷo āpyo.}
\end{quote}

\begin{quote}
Every family contributed to the temple.
\end{quote}

A masculine noun.

\begin{cousinweb}
\textbf{\gu{ફાળો}} (\emph{phāḷo}) CDIAL lists Gujarati \emph{phāḷo}, \textquotedblleft{}a share\textquotedblright{}, under a reconstructed \emph{phāla}, \textquotedblleft{}a piece split off\textquotedblright{} (CDIAL 9073): one's share of a common cost.
\end{cousinweb}

\begin{grammarlens}[title={What you've built}]
One more word for giving reasons and opinions.
\end{grammarlens}

\subsection*{Your turn}
\begin{itemize}
\raggedright
  \item \emph{Say it:} \emph{phāḷo}
  \item \emph{Say it:} \emph{phāḷo}, once more
  \item \emph{Say it:} \emph{mandir māṭe darek kuṭumbe phāḷo āpyo.} — Every family contributed to the temple.
  \item \emph{Recall:} say the Gujarati for to appoint, then the Gujarati for a tax, then say \emph{phāḷo} again
\end{itemize}

\begin{tcolorbox}[breakable,colback=teal!4,colframe=teal!35!black,title={Before you move on}]
What does \gu{ફાળો} mean? (\textquotedblleft{}A contribution\textquotedblright{}.) Say it once more.
\end{tcolorbox}
\section[\texorpdfstring{\gu{આર્થિક} (ārthik)}{આર્થિક (ārthik)}]{\gu{આર્થિક} (ārthik) — economic; financial}

\label{lesson:GU-C355-arthik}



\begin{quote}
Before the new one: say the Gujarati for a tax, then the Gujarati for a contribution.
\end{quote}

\subsection*{You'll want to know: \gu{આર્થિક}}
\textbf{\gu{આર્થિક}} — \emph{ārthik} — \textquotedblleft{}economic; financial\textquotedblright{}.

\begin{quote}
\gu{નોકરી} \gu{ગયા} \gu{પછી} \gu{અરુણની} \gu{આર્થિક} \gu{પરિસ્થિતિ} \gu{નબળી} \gu{થઈ}.
\end{quote}

\begin{quote}
\emph{nokrī gayā pachī Aruṇnī ārthik paristhiti nabḷī thaī.}
\end{quote}

\begin{quote}
After he lost his job, Arun's finances grew weak.
\end{quote}

\begin{cousinweb}
\textbf{\gu{આર્થિક}} (\emph{ārthik}) is a tatsama, borrowed from Sanskrit \emph{ārthika}, built on \emph{artha}, \textquotedblleft{}aim, wealth\textquotedblright{} (Monier-Williams).
\end{cousinweb}

\begin{grammarlens}[title={What you've built}]
One more word for giving reasons and opinions.
\end{grammarlens}

\subsection*{Your turn}
\begin{itemize}
\raggedright
  \item \emph{Say it:} \emph{ārthik}
  \item \emph{Say it:} \emph{ārthik}, once more
  \item \emph{Say it:} \emph{nokrī gayā pachī Aruṇnī ārthik paristhiti nabḷī thaī.} — After he lost his job, Arun's finances grew weak.
  \item \emph{Recall:} say the Gujarati for a tax, then the Gujarati for a contribution, then say \emph{ārthik} again
\end{itemize}

\begin{tcolorbox}[breakable,colback=teal!4,colframe=teal!35!black,title={Before you move on}]
What does \gu{આર્થિક} mean? (\textquotedblleft{}Economic; financial\textquotedblright{}.) Say it once more.
\end{tcolorbox}
\section[\texorpdfstring{\gu{ગ્રાન્ટ} (grānṭ)}{ગ્રાન્ટ (grānṭ)}]{\gu{ગ્રાન્ટ} (grānṭ) — a grant (of public money)}

\label{lesson:GU-C355-grant}



\begin{quote}
Before the new one: say the Gujarati for a contribution, then the Gujarati for economic.
\end{quote}

\subsection*{You'll want to know: \gu{ગ્રાન્ટ}}
\textbf{\gu{ગ્રાન્ટ}} — \emph{grānṭ} — \textquotedblleft{}a grant (of public money)\textquotedblright{}.

\begin{quote}
\gu{સરકારે} \gu{શાળાને} \gu{નવી} \gu{ગ્રાન્ટ} \gu{આપી}.
\end{quote}

\begin{quote}
\emph{sarkāre śāḷāne navī grānṭ āpī.}
\end{quote}

\begin{quote}
The government gave the school a new grant.
\end{quote}

A feminine noun.

\begin{cousinweb}
\textbf{\gu{ગ્રાન્ટ}} (\emph{grānṭ}) is borrowed from English \emph{grant}.
\end{cousinweb}

\begin{grammarlens}[title={What you've built}]
One more word for giving reasons and opinions.
\end{grammarlens}

\subsection*{Your turn}
\begin{itemize}
\raggedright
  \item \emph{Say it:} \emph{grānṭ}
  \item \emph{Say it:} \emph{grānṭ}, once more
  \item \emph{Say it:} \emph{sarkāre śāḷāne navī grānṭ āpī.} — The government gave the school a new grant.
  \item \emph{Recall:} say the Gujarati for a contribution, then the Gujarati for economic, then say \emph{grānṭ} again
\end{itemize}

\begin{tcolorbox}[breakable,colback=teal!4,colframe=teal!35!black,title={Before you move on}]
What does \gu{ગ્રાન્ટ} mean? (\textquotedblleft{}A grant (of public money)\textquotedblright{}.) Say it once more.
\end{tcolorbox}
\section[\texorpdfstring{\gu{ઉઘરાવવું} (ughrāvvũ)}{ઉઘરાવવું (ughrāvvũ)}]{\gu{ઉઘરાવવું} (ughrāvvũ) — to collect money (from people)}

\label{lesson:GU-C355-ughravvu}



\begin{quote}
Before the new one: say the Gujarati for economic, then the Gujarati for a grant.
\end{quote}

\subsection*{You'll want to know: \gu{ઉઘરાવવું}}
\textbf{\gu{ઉઘરાવવું}} — \emph{ughrāvvũ} — \textquotedblleft{}to collect money (from people)\textquotedblright{}.

\begin{quote}
\gu{ગામના} \gu{જુવાનો} \gu{ઉજાણી} \gu{માટે} \gu{ફાળો} \gu{ઉઘરાવે} \gu{છે}.
\end{quote}

\begin{quote}
\emph{gāmnā juvāno ujāṇī māṭe phāḷo ughrāve che.}
\end{quote}

\begin{quote}
The young people of the village are collecting money for a picnic.
\end{quote}

Past, with \gu{તેણે} (he or she …): \gu{તેણે} \gu{બધો} \gu{ફાળો} \gu{ઉઘરાવ્યો}, \emph{teṇe badho phāḷo ughrāvyo}.

\begin{cousinweb}
\textbf{\gu{ઉઘરાવવું}} (\emph{ughrāvvũ}) is inherited from \emph{udgharati}, a form CDIAL reconstructs; CDIAL lists the Gujarati verb, \textquotedblleft{}to collect money promised\textquotedblright{} (CDIAL 1971).
\end{cousinweb}

\begin{grammarlens}[title={What you've built}]
One more word for giving reasons and opinions.
\end{grammarlens}

\subsection*{Your turn}
\begin{itemize}
\raggedright
  \item \emph{Say it:} \emph{ughrāvvũ}
  \item \emph{Say it:} \emph{ughrāvvũ}, once more
  \item \emph{Say it:} \emph{gāmnā juvāno ujāṇī māṭe phāḷo ughrāve che.} — The young people of the village are collecting money for a picnic.
  \item \emph{Recall:} say the Gujarati for economic, then the Gujarati for a grant, then say \emph{ughrāvvũ} again
\end{itemize}

\begin{tcolorbox}[breakable,colback=teal!4,colframe=teal!35!black,title={Before you move on}]
What does \gu{ઉઘરાવવું} mean? (\textquotedblleft{}To collect money (from people)\textquotedblright{}.) Say it once more.
\end{tcolorbox}
